% GENERATED FILE. Edit canonical lessons, then run npm run generate:books.
% canonical-source-hash: fnv1a64:829a44ebc523b3d9
% canonical-lessons: ZH-C220-jianshen, ZH-C220-jianfei, ZH-C220-chuhan, ZH-C220-tiaosheng, ZH-C220-zuocao

\chapter{Keeping Fit}
\label{ch:zh-keeping-fit}
\hlchaptermodality{220}

\begin{chapteropening}
\textbf{By the end of this chapter:} \emph{I can talk about working out, losing weight, sweating, skipping and doing exercises.}

Complete the last lesson of chapter 220: I can talk about working out, losing weight, sweating, skipping and doing exercises.
\end{chapteropening}

\section[jiànshēn]{jiànshēn — to work out; to keep fit}

\label{lesson:ZH-C220-jianshen}



\begin{quote}
Before the new one: say the Mandarin for air temperature, then the Mandarin for to clear up.
\end{quote}

\subsection*{You'll want to know: jiànshēn}
\emph{jiànshēn} — \textquotedblleft{}to work out; to keep fit\textquotedblright{}.

\begin{quote}
\emph{wǒ měitiān xiàwǔ qù jiànshēn} — I go and work out every afternoon.
\end{quote}

Its characters come later; for now it is a word for the ear.

\begin{cousinweb}
\emph{jiànshēn} is built from \emph{jiàn} 'strong' and \emph{shēn} 'body', both inherited from Middle Chinese, \emph{gjonH} and \emph{syin} (Guangyun, in Baxter's transcription), which glosses the first 'sturdy' and the second 'oneself' and 'the body'.
\end{cousinweb}

\begin{grammarlens}[title={What you've built}]
One more word for saying what I do.
\end{grammarlens}

\subsection*{Your turn}
\begin{itemize}
\raggedright
  \item \emph{Say it:} \emph{jiànshēn}
  \item \emph{Say it:} \emph{jiànshēn}, once more
  \item \emph{Say it:} \emph{wǒ měitiān xiàwǔ qù jiànshēn} — I go and work out every afternoon.
  \item \emph{Recall:} say the Mandarin for air temperature, then the Mandarin for to clear up, then say \emph{jiànshēn} again
\end{itemize}

\begin{tcolorbox}[breakable,colback=teal!4,colframe=teal!35!black,title={Before you move on}]
What does jiànshēn mean? (\textquotedblleft{}To work out; to keep fit\textquotedblright{}.) Say it once more.
\end{tcolorbox}
\section[jiǎnféi]{jiǎnféi — to lose weight; to diet}

\label{lesson:ZH-C220-jianfei}



\begin{quote}
Before the new one: say the Mandarin for to clear up, then the Mandarin for to work out.
\end{quote}

\subsection*{You'll want to know: jiǎnféi}
\emph{jiǎnféi} — \textquotedblleft{}to lose weight; to diet\textquotedblright{}.

\begin{quote}
\emph{wǒ jiějie zhèngzài jiǎnféi} — My big sister is trying to lose weight.
\end{quote}

Its characters come later; for now it is a word for the ear.

\begin{cousinweb}
\emph{jiǎnféi} is built from \emph{jiǎn} 'to reduce' and \emph{féi} 'fat'. The second is inherited from Middle Chinese \emph{bj+j} (Guangyun, in Baxter's transcription), explained after an older book as having much flesh; the dictionary glosses the first, \emph{keamX} there, 'to lessen'.
\end{cousinweb}

\begin{grammarlens}[title={What you've built}]
One more word for saying what I do.
\end{grammarlens}

\subsection*{Your turn}
\begin{itemize}
\raggedright
  \item \emph{Say it:} \emph{jiǎnféi}
  \item \emph{Say it:} \emph{jiǎnféi}, once more
  \item \emph{Say it:} \emph{wǒ jiějie zhèngzài jiǎnféi} — My big sister is trying to lose weight.
  \item \emph{Recall:} say the Mandarin for to clear up, then the Mandarin for to work out, then say \emph{jiǎnféi} again
\end{itemize}

\begin{tcolorbox}[breakable,colback=teal!4,colframe=teal!35!black,title={Before you move on}]
What does jiǎnféi mean? (\textquotedblleft{}To lose weight; to diet\textquotedblright{}.) Say it once more.
\end{tcolorbox}
\section[chūhàn]{chūhàn — to sweat}

\label{lesson:ZH-C220-chuhan}



\begin{quote}
Before the new one: say the Mandarin for to work out, then the Mandarin for to lose weight.
\end{quote}

\subsection*{You'll want to know: chūhàn}
\emph{chūhàn} — \textquotedblleft{}to sweat\textquotedblright{}.

\begin{quote}
\emph{tiān tài rè le wǒ chūhàn le} — It's so hot that I'm sweating.
\end{quote}

Its characters come later; for now it is a word for the ear.

\begin{cousinweb}
\emph{chūhàn} is built from \zh{出} 'to come out' and \emph{hàn} 'sweat', both inherited from Middle Chinese, \emph{tsyhwit} and \emph{hanH} (Guangyun, in Baxter's transcription), which glosses the second 'hot sweat'.
\end{cousinweb}

\begin{grammarlens}[title={What you've built}]
One more word for saying what I do.
\end{grammarlens}

\subsection*{Your turn}
\begin{itemize}
\raggedright
  \item \emph{Say it:} \emph{chūhàn}
  \item \emph{Say it:} \emph{chūhàn}, once more
  \item \emph{Say it:} \emph{tiān tài rè le wǒ chūhàn le} — It's so hot that I'm sweating.
  \item \emph{Recall:} say the Mandarin for to work out, then the Mandarin for to lose weight, then say \emph{chūhàn} again
\end{itemize}

\begin{tcolorbox}[breakable,colback=teal!4,colframe=teal!35!black,title={Before you move on}]
What does chūhàn mean? (\textquotedblleft{}To sweat\textquotedblright{}.) Say it once more.
\end{tcolorbox}
\section[tiàoshéng]{tiàoshéng — to skip with a rope}

\label{lesson:ZH-C220-tiaosheng}



\begin{quote}
Before the new one: say the Mandarin for to lose weight, then the Mandarin for to sweat.
\end{quote}

\subsection*{You'll want to know: tiàoshéng}
\emph{tiàoshéng} — \textquotedblleft{}to skip with a rope\textquotedblright{}.

\begin{quote}
\emph{wǒ nǚ'ér měitiān tiàoshéng} — My daughter skips every day.
\end{quote}

Its characters come later; for now it is a word for the ear.

\begin{cousinweb}
\emph{tiàoshéng} is built from \emph{tiào} 'to jump' and \emph{shéng} 'rope'. The first is inherited from Middle Chinese \emph{dew} (Guangyun, in Baxter's transcription), glossed 'to leap', whose regular outcome is tiáo, not the falling tiào; the dictionary glosses the second, \emph{zying} there, 'straight' and 'a cord'.
\end{cousinweb}

\begin{grammarlens}[title={What you've built}]
One more word for saying what I do.
\end{grammarlens}

\subsection*{Your turn}
\begin{itemize}
\raggedright
  \item \emph{Say it:} \emph{tiàoshéng}
  \item \emph{Say it:} \emph{tiàoshéng}, once more
  \item \emph{Say it:} \emph{wǒ nǚ'ér měitiān tiàoshéng} — My daughter skips every day.
  \item \emph{Recall:} say the Mandarin for to lose weight, then the Mandarin for to sweat, then say \emph{tiàoshéng} again
\end{itemize}

\begin{tcolorbox}[breakable,colback=teal!4,colframe=teal!35!black,title={Before you move on}]
What does tiàoshéng mean? (\textquotedblleft{}To skip with a rope\textquotedblright{}.) Say it once more.
\end{tcolorbox}
\section[zuòcāo]{zuòcāo — to do a fitness routine; to do physical exercises}

\label{lesson:ZH-C220-zuocao}



\begin{quote}
Before the new one: say the Mandarin for to sweat, then the Mandarin for to skip with a rope.
\end{quote}

\subsection*{You'll want to know: zuòcāo}
\emph{zuòcāo} — \textquotedblleft{}to do a fitness routine; to do physical exercises\textquotedblright{}.

\begin{quote}
\emph{wǒ bàba měitiān zǎoshang zuòcāo} — My dad does his exercises every morning.
\end{quote}

Its characters come later; for now it is a word for the ear.

\begin{cousinweb}
\emph{zuòcāo} is built from \zh{做} 'to do' and \emph{cāo} 'a drill'. The second is inherited from Middle Chinese \emph{tshaw} (Guangyun, in Baxter's transcription), glossed 'to hold'; how it came to mean a drill is not settled. No dictionary we use records an older form of the first, so its history is not settled.
\end{cousinweb}

\begin{grammarlens}[title={What you've built}]
One more word for saying what I do.
\end{grammarlens}

\subsection*{Your turn}
\begin{itemize}
\raggedright
  \item \emph{Say it:} \emph{zuòcāo}
  \item \emph{Say it:} \emph{zuòcāo}, once more
  \item \emph{Say it:} \emph{wǒ bàba měitiān zǎoshang zuòcāo} — My dad does his exercises every morning.
  \item \emph{Recall:} say the Mandarin for to sweat, then the Mandarin for to skip with a rope, then say \emph{zuòcāo} again
\end{itemize}

\begin{tcolorbox}[breakable,colback=teal!4,colframe=teal!35!black,title={Before you move on}]
What does zuòcāo mean? (\textquotedblleft{}To do a fitness routine; to do physical exercises\textquotedblright{}.) Say it once more.
\end{tcolorbox}
